\documentclass[11pt]{article}
\usepackage[margin=25mm]{geometry}
\usepackage{fontspec}\setmainfont{Latin Modern Roman}
\usepackage{amsmath,amssymb,mathtools}
\usepackage{unicode-math}\setmathfont{Latin Modern Math}
\usepackage{xurl,hyperref}
\hypersetup{colorlinks=true,urlcolor=blue}
\setcounter{secnumdepth}{0}
\setlength{\parindent}{0pt}\setlength{\parskip}{5pt}
\setlength{\emergencystretch}{3em}
\newcommand{\sref}[2]{\hyperlink{source:#1:#2}{[S#2]}}
\newcommand{\cataloguescope}[1]{\paragraph{Recorded target.}#1}
\begin{document}
\section{Kirby Problem 1.22}\label{statement}

\subsection{Definitions and mathematical statement}

\paragraph{Shared notation.}$\mathbb N=\{1,2,\ldots\}$, and $\mathbb Z,\mathbb Q,\mathbb R,\mathbb C$ denote the integers, rationals, reals and complex numbers. Any other conventions are specified locally.


Use the integral hat Heegaard Floer group $\widehat{\mathrm{HF}}(Y)$ and integral hat knot Floer group $\widehat{\mathrm{HFK}}(Y,K)$, with rank meaning the rank of the free abelian part. Knot Floer isomorphisms preserve Maslov and Alexander gradings. A rational homology sphere $Y$ is an L-space when $\operatorname{rank}\widehat{\mathrm{HF}}(Y)=|H_1(Y;\mathbb Z)|$. An L-space knot is a knot in $S^3$ having a positive integral surgery that is an L-space. A hyperbolic knot has a complete finite-volume hyperbolic complement. Write $\Sigma_2(K)$ for the two-fold cover of $S^3$ branched along $K$.
\paragraph{Statement recorded in the supplied source.}
For every hyperbolic L-space knot $K$, is $\Sigma_2(K)$ not an L-space? \sref{statement}{1}\sref{statement}{2}
\subsection{Short English statement}
The double branched cover of a hyperbolic L-space knot should fail to be an L-space.
\subsection{Sources}
\begin{description}
\item[\hypertarget{source:statement:1}{S1}] ULAM.AI and the UnsolvedMath Contributors, \emph{UnsolvedMath: A Curated Collection of Open Mathematics Problems}, catalogue entry. CC BY 4.0. \url{https://huggingface.co/datasets/ulamai/UnsolvedMath}.
\item[\hypertarget{source:statement:2}{S2}] R. \.{I}nan\c{c} Baykur, Robion C. Kirby and Daniel Ruberman, K3 -- A New Problem List in Low-Dimensional Topology, author preliminary edition (2026), Problem 1.22 and its remarks. \url{https://math.berkeley.edu/sites/default/files/surv-295-ruberman-watermarked-author-pdf.pdf}.
\end{description}
\label{end:statement}
\end{document}
